% Remove the following line and define them in main.tex
\keywords{corporate disclosure analytics; OCR pipeline; tone classification; ontology mapping; greenwashing screening}

\begin{thesisabstract}{english}
  \abstractmetadata

This thesis develops an executable framework for analysing Indonesian
sustainability reports as structured, record-level ESG evidence. The workflow
combines OCR-based document expansion, page-level provenance tracking, large
language model extraction, record-level schema design, ontology alignment, and
comparison with ClimateBERT-style climate labels.

The framework was evaluated on an active thesis subset of 23 processed reports
covering approximately 5,512 pages. It produced 332 tone-bearing ESG records
and 2,074 downstream analytical rows, while all 52 tracked aspects were mapped
to ontology paths. The results show that disclosure tone is measurable and
useful: commitment was the dominant tone, environmental content was the
dominant pillar, and tone labels reached 83.7\% agreement with
ClimateBERT-style commitment labels with Cohen's kappa of 0.645.

Extraction quality nevertheless depended strongly on prompt design and model
compliance, with missing-tone cases and schema drift as the main weaknesses.
The study concludes that tone-aware, provenance-preserving ESG disclosure
analysis is feasible and informative, but stronger expert annotation and
broader validation are still required before benchmark-level claims can be
made.
\end{thesisabstract}
